\chapter{Couverture d'un portefeuille de produits structurés}\label{ch:s2:hedging-a-structured-products-book}

Les particuliers achètent des autocallables en gros volumes. Les banques qui les vendent se retrouvent longues de volatilité qu'elles doivent vendre, courtes de corrélation et longues de
dividendes, et leur couverture fait bouger ces marchés. La Banque des règlements internationaux a relevé en septembre 2026 que la couverture des autocallables largement
vendus aux particuliers coréens peut exiger des transactions importantes et déstabilisantes près des barrières. Sur le book synthétique de ce chapitre, portant sur 1~milliard d'euros de
d'autocallables à cinq ans, la banque est longue de 9.66~millions d'euros de vega par point de volatilité et dégage une marge de 8.31~millions d'euros. Vendre ce vega
sur un marché qui absorbe 10~millions d'euros par point de concession abaisse la volatilité implicite à long terme de 0.97 point et coûte à la banque 4.67~millions d'euros,
soit plus de la moitié de sa marge. La construction est \texttt{firm.sphedge}.

\section{Les expositions du book d'autocallables}

\begin{definition}[Exposition du portefeuille d'exotiques]\label{def:s2:hedging-a-structured-products-book:exposure}
Une \emph{exposition du portefeuille d'exotiques}\index{exposition du portefeuille d'exotiques} est une sensibilité d'un book de produits structurés à une donnée de marché que les instruments vanille
ne couvrent qu'en partie (volatilité aux longues échéances, skew, dividendes, corrélation), mesurée en réévaluant l'ensemble du book
la donnée étant choquée.
\end{definition}

La banque a vendu trois notes (\cref{lst:s2:hedging-a-structured-products-book:book})~: 400~millions d'euros sur l'indice A seul, 300~millions d'euros sur l'indice B
seul, et 300~millions d'euros sur le pire des deux. Chacune est l'autocallable du livre~5 (chapitre~18)~: cinq ans, observations annuelles, rappel automatique à 100~\%, un
coupon Phoenix de 7~\% (9~\% pour le pire des deux) versé au-dessus d'une barrière de 70~\% avec mémoire, et capital protégé sauf si l'indice termine sous 60~\%. Les
indices ont une volatilité de 20~\%, un skew de volatilité locale qui ajoute 10 points par baisse de 100~\%, des dividendes de 3~\% et une corrélation de 0.7. Les notes isolées
valent 99.12 pour 100, et le pire des deux 99.29~: la marge de la banque est de 8.31~millions d'euros.

\begin{center}
\small
\begin{tabular}{@{}lcccc@{}}
\toprule
P\&L de l'émetteur pour le choc (millions d'euros) & A seul & B seul & pire des deux & book\\
\midrule
volatilité $+1$ point & 3.22 & 2.42 & 4.02 & 9.66\\
skew $+0.05$ & 1.47 & 1.10 & 1.86 & 4.42\\
rendement de dividende $+0.1$ point & 0.51 & 0.38 & 0.55 & 1.43\\
corrélation $+0.05$ & & & $-1.05$ & $-1.05$\\
\bottomrule
\end{tabular}
\end{center}

La note de l'investisseur est, pour une large part, un put vendu qui se déclenche sous 60~\%. La banque détient l'autre côté~: elle gagne lorsque la volatilité, le skew à la baisse ou
les dividendes montent, ce qui renchérit le put. La note sur le pire des deux ajoute une position courte en corrélation~: quand les indices évoluent plus de concert, le pire
des deux a moins de chances d'être mauvais, ce qui aide l'investisseur et coûte à la banque. Ces expositions sont l'image miroir de ce que l'investisseur a vendu, et la
banque doit s'en décharger.

\begin{dated}{2026-09}{Autocallables en Corée}\label{dat:s2:hedging-a-structured-products-book:korea}
La BIS Quarterly Review de septembre 2026 décrit les autocallables comme largement vendus aux particuliers coréens et note que leur couverture peut exiger
des transactions importantes et déstabilisantes près des barrières. Les émissions référencées sur Samsung Electronics et SK Hynix, qui pèsent ensemble plus de la moitié du KOSPI 200, ont
fortement progressé au cours du premier semestre 2026.
\end{dated}

\section{Vega et skew}

Les expositions du book évoluent avec l'indice (\cref{fig:s2:hedging-a-structured-products-book:profile}). À l'origine, la note isolée sur A laisse la banque
longue de 3.22~millions d'euros de vega, et sa couverture détient 106~millions d'euros d'indice. Si l'indice ouvrait plus bas, le vega monterait jusqu'à un maximum de 4.45
millions d'euros vers 80~\% du prix d'exercice, puis baisserait et deviendrait négatif sous environ 52~\%~: une fois la protection probablement perdue, la note est surtout une détention de
l'indice. La couverture augmente à mesure que l'indice baisse, jusqu'à 334~millions d'euros à 55~\%, ce qui signifie acheter dans un marché en baisse. En dessous, elle se réduit, jusqu'à
246~millions d'euros à 40~\%, ce qui signifie vendre à mesure que le marché poursuit sa chute. Cette seconde zone est la transaction déstabilisante près de la barrière.
(\cref{lst:s2:hedging-a-structured-products-book:profile} calcule les deux courbes.)

\begin{omfigure}
\begin{tikzpicture}
\begin{axis}[width=11cm,height=5.4cm,axis y line*=right,axis x line=none,ylabel={\small indice détenu en couverture (\euro{} M)},
  tick label style={font=\footnotesize},xmin=40,xmax=120,ymin=0,ymax=400,table/col sep=comma]
\addplot[omDef,thick,dashed] table[x=spot,y=hedge]{figdata/strategies-2/28-hedging-a-structured-products-book/profile.csv};\label{plot:s2:sphedge}
\end{axis}
\begin{axis}[width=11cm,height=5.4cm,axis y line*=left,xlabel={\small niveau de l'indice (\% du prix d'exercice)},ylabel={\small vega de l'émetteur (\euro{} M par point)},
  tick label style={font=\footnotesize},xmin=40,xmax=120,ymin=-3,ymax=5,grid=major,grid style={gray!25},table/col sep=comma,
  legend style={font=\scriptsize,at={(0.5,-0.3)},anchor=north},legend columns=2]
\addplot[omThm,thick,mark=*,mark size=1pt] table[x=spot,y=vega]{figdata/strategies-2/28-hedging-a-structured-products-book/profile.csv};
\addlegendentry{vega (gauche)}
\addlegendimage{/pgfplots/refstyle=plot:s2:sphedge}\addlegendentry{indice détenu en couverture delta (droite)}
\end{axis}
\end{tikzpicture}

\omcaption{La couverture en vega et en delta de la banque pour la note isolée de 400~millions d'euros sur l'indice A, à l'origine, si l'indice se situait à chaque niveau par rapport
au prix d'exercice. Données~: \texttt{s2\_sphedge.profile}.}\label{fig:s2:hedging-a-structured-products-book:profile}
\end{omfigure}

\section{Corrélation et dividendes}

\begin{definition}[Trade de recyclage]\label{def:s2:hedging-a-structured-products-book:recycling}
Un \emph{trade de recyclage}\index{trade de recyclage} est une transaction par laquelle un émetteur de produits structurés transfère à d'autres investisseurs les expositions que son
book ne peut pas couvrir avec des instruments vanille~: vente de volatilité longue, vente de dividendes, achat de corrélation, généralement moyennant une concession.
\end{definition}

L'exposition de la banque aux dividendes, 1.43~millions d'euros par 0.1 point de rendement de dividende, est recyclée par la vente de contrats à terme sur dividendes ou de swaps, l'offre que le chapitre~26
a modélisée~; Eurex décrit ses dérivés sur dividendes comme servant à couvrir les dividendes, notamment pour les produits structurés. La corrélation courte, 1.05~millions d'euros
par 0.05 de corrélation, est recyclée par l'achat de corrélation~: la vente de volatilité de l'indice contre la volatilité des titres individuels, l'autre face du
trade de dispersion du chapitre~2. Chaque \omterm{def:s2:hedging-a-structured-products-book:recycling}{trade de recyclage} a besoin d'un acheteur, et les acheteurs (hedge funds, traders de dispersion, investisseurs en dividendes) sont rémunérés pour
prendre l'autre côté (\cref{fig:s2:hedging-a-structured-products-book:exposures}).

\begin{omfigure}
\begin{tikzpicture}
\begin{axis}[width=11cm,height=5.2cm,ybar,bar width=8pt,ylabel={\small P\&L de l'émetteur pour le choc (\euro{} M)},tick label style={font=\footnotesize},
  xtick=data,xticklabels={vega,skew,dividendes,corrélation},ymin=-1.5,ymax=4.5,grid=major,grid style={gray!25},table/col sep=comma,
  legend style={font=\scriptsize,at={(0.97,0.97)},anchor=north east}]
\addplot[fill=omThm!60,draw=omThm] table[x=k,y=n0]{figdata/strategies-2/28-hedging-a-structured-products-book/exposures.csv};
\addplot[fill=gray!40,draw=gray] table[x=k,y=n1]{figdata/strategies-2/28-hedging-a-structured-products-book/exposures.csv};
\addplot[fill=omDef!60,draw=omDef] table[x=k,y=n2]{figdata/strategies-2/28-hedging-a-structured-products-book/exposures.csv};
\legend{A seul,B seul,pire des deux}
\end{axis}
\end{tikzpicture}

\omcaption{Les expositions du book synthétique par note~: le P\&L de la banque pour une hausse de 1 point de volatilité, une hausse de 0.05 du skew, une hausse de 0.1 point du rendement de
dividende et une hausse de 0.05 de la corrélation. Données~: \texttt{s2\_sphedge.table}.}\label{fig:s2:hedging-a-structured-products-book:exposures}
\end{omfigure}

\section{Trades de recyclage}

Vendre 9.66~millions d'euros de vega par point fait bouger le marché qui l'absorbe. Si les acheteurs prennent $D$~millions d'euros de vega pour chaque point de concession, la volatilité implicite à long terme
baisse de $9.66/D$ points~; la banque paie en moyenne la moitié de ce mouvement, et un acheteur qui conserve sa position jusqu'au retour de la volatilité en encaisse la totalité~:

\begin{center}
\small
\begin{tabular}{@{}lccc@{}}
\toprule
profondeur (millions d'euros de vega par point) & 5 & 10 & 20\\
\midrule
mouvement de la volatilité implicite à long terme (points) & $-1.93$ & $-0.97$ & $-0.48$\\
coût pour la banque (millions d'euros) & 9.33 & 4.67 & 2.33\\
gain de l'acheteur si la volatilité revient (millions d'euros) & 18.67 & 9.33 & 4.67\\
\bottomrule
\end{tabular}
\end{center}

Le coût du recyclage est de l'ordre de la marge, soit 8.31~millions d'euros. Lorsque de nombreux émetteurs vendent les mêmes expositions, la profondeur se partage et le mouvement
devient plus large. L'offre qui écarte de sa valeur de marché la volatilité à long terme, les dividendes et la corrélation est précisément ce que rémunèrent les acheteurs des chapitres~2, 3 et~26.
Les alternatives de l'émetteur sont de garder le risque, ce que ses limites de risque autorisent rarement, ou d'intégrer le coût du recyclage dans le prix des notes.

\section{Fiches de stratégie}

\begin{strategyfile}{Recyclage du vega vers les hedge funds}\label{strat:s2:hedging-a-structured-products-book:vega}
\sfield{Qui vous paie, et pourquoi} L'émetteur, qui doit vendre le vega à long terme que lui laissent ses notes.
\sfield{Instruments et plateformes} Options sur indices à long terme et swaps de variance, de gré à gré.
\sfield{Signal} Volatilité implicite à long terme inférieure à son niveau d'équilibre après les vagues d'émissions.
\sfield{Dimensionnement et exécution} Acheter ce que vendent les émetteurs~; conserver la position jusqu'au retour de la volatilité.
\sfield{Coûts} Fourchettes larges sur les options à long terme~; financement de la prime.
\sfield{Comment elle meurt} Des émissions qui ne cessent jamais~; un régime de volatilité qui reste bas.
\sfield{Horizon, capacité, infrastructure} Des années.
\sfield{Backtest honnête} Historique de la volatilité implicite à long terme tel qu'il a été coté, avec ses fourchettes.
\sfield{Sources} Ce chapitre~: 0.97 point de volatilité implicite pour 9.66~millions d'euros de vega à une profondeur de 10~millions d'euros.
\end{strategyfile}

\begin{strategyfile}{Recyclage des dividendes}\label{strat:s2:hedging-a-structured-products-book:dividends}
\sfield{Qui vous paie, et pourquoi} Les émetteurs, longs de dividendes au travers des couvertures de leurs notes.
\sfield{Instruments et plateformes} Contrats à terme et swaps sur dividendes d'indices.
\sfield{Signal} Contrats à terme sur dividendes inférieurs aux dividendes prévus.
\sfield{Dimensionnement et exécution} Conservée jusqu'à l'échéance.
\sfield{Coûts} Fourchettes des contrats à terme~; baisses liées à la récession.
\sfield{Comment elle meurt} Les récessions et les interdictions de distribution de dividendes.
\sfield{Horizon, capacité, infrastructure} De un à cinq ans.
\sfield{Backtest honnête} Inclure les années de récession.
\sfield{Sources} Fiche Eurex~; chapitre~26~; ce chapitre~: 1.43~million d'euros par 0.1 point de rendement de dividende.
\end{strategyfile}

\begin{strategyfile}{Recyclage de la corrélation}\label{strat:s2:hedging-a-structured-products-book:corr}
\sfield{Qui vous paie, et pourquoi} Les émetteurs de notes worst-of, vendeurs de corrélation.
\sfield{Instruments et plateformes} Swaps de corrélation~; indice contre volatilité des valeurs individuelles (dispersion).
\sfield{Signal} Corrélation implicite comparée à la \omterm{def:s2:dispersion-and-correlation:realised}{corrélation réalisée}.
\sfield{Dimensionnement et exécution} Dispersion pondérée par le vega (chapitre~2).
\sfield{Coûts} Fourchettes des options sur valeurs individuelles.
\sfield{Comment elle meurt} Les krachs, lorsque la corrélation fait un saut.
\sfield{Horizon, capacité, infrastructure} De quelques mois à quelques années.
\sfield{Backtest honnête} Corrélation pendant les krachs.
\sfield{Sources} Ce chapitre~: $-$1.05~million d'euros par 0.05 de corrélation sur un worst-of de 300~millions d'euros.
\end{strategyfile}

\begin{strategyfile}{Programme de couverture du skew}\label{strat:s2:hedging-a-structured-products-book:skew}
\sfield{Qui vous paie, et pourquoi} Les limites de risque propres à l'émetteur~; le prix que le marché accorde à la protection contre la baisse.
\sfield{Instruments et plateformes} Puts sur indices hors de la monnaie et put spreads.
\sfield{Signal} L'exposition au skew du portefeuille par échéance et par niveau.
\sfield{Dimensionnement et exécution} Vendre le skew à la baisse contre la position longue du portefeuille, davantage à mesure que l'indice approche des barrières.
\sfield{Coûts} Put spreads~; la couverture déstabilisante près des barrières.
\sfield{Comment elle meurt} Un krach qui franchit les barrières.
\sfield{Horizon, capacité, infrastructure} La durée de vie des notes.
\sfield{Backtest honnête} Trajectoires qui traversent les barrières, et pas seulement les trajectoires calmes.
\sfield{Sources} BIS Quarterly Review, septembre 2026~; ce chapitre~: 4.42~millions d'euros par 0.05 de skew.
\end{strategyfile}

\section{Tutoriel~: acheter ce que personne ne veut}\label{tut:s2:hedging-a-structured-products-book:tut}

\begin{tutorial}
\textbf{Objectif.} Valoriser un portefeuille d'autocallables, mesurer ses expositions par revalorisation, suivre le vega d'une note et la couvrir contre le niveau de l'indice, et chiffrer le
recyclage de son vega. \textbf{État final~:} les trois tableaux et les deux figures.
\begin{tutsteps}
\item \textbf{Le portefeuille et ses expositions}.
\omcode{code/firm/sphedge/firm_sphedge.py}{56}{79}{Notes sur les trajectoires mensuelles~; le P\&L de l'émetteur pour chaque choc.}\label{lst:s2:hedging-a-structured-products-book:book}
\item \textbf{Au travers du niveau de l'indice, et recyclage}.
\omcode{code/firm/sphedge/firm_sphedge.py}{82}{102}{Vega et delta par niveau~; le mouvement et le coût de la vente de vega.}\label{lst:s2:hedging-a-structured-products-book:profile}
\item \textbf{Exécuter} \texttt{table()}, \texttt{margin()}, \texttt{profile()}, \texttt{recycling()} et \texttt{fig\_sphedge.py}.
\end{tutsteps}
\textbf{Pour aller plus loin.} Vieillir le portefeuille d'un an et recalculer~; ajouter un troisième indice et un worst-of sur trois~; laisser la profondeur baisser lorsque plusieurs émetteurs vendent
en même temps.
\end{tutorial}

\section{Construction~: couverture de produits structurés}\label{bld:s2:hedging-a-structured-products-book:sphedge}

\begin{build}
\bfield{Rôle} Expositions d'un portefeuille d'autocallables par revalorisation, leur profil selon le niveau de l'indice, et le coût du recyclage.
\bfield{Interface} \texttt{Note}, \texttt{BOOK}, \texttt{MarketState(\dots)}, \texttt{note\_value(note, mkt, paths, seed, spot)},
\texttt{book\_exposures}, \texttt{spot\_profile}, \texttt{recycling\_impact(vega, depth)}.
\bfield{Règles} Termes et flux de trésorerie du livre~5~; nombres aléatoires communs pour chaque choc~; pas mensuels (observations annuelles, knock-in à l'échéance).
\bfield{Tests d'acceptation} \texttt{code/firm/sphedge/tests/}~: les signes du portefeuille~; des valeurs de note cohérentes~; le gain du worst-of dû à la corrélation~; le recyclage refait à la main.
by hand.
\bfield{Pour aller plus loin} Knock-in quotidien~; vieillissement~; plusieurs émetteurs partageant la profondeur.
\end{build}

\begin{omsources}
\item Bank for International Settlements, \emph{BIS Quarterly Review}, September 2026, Box B, ``Short on memory? Leveraged ETFs and chip stock
volatility''.
\item Eurex, ``Dividend derivatives'', factsheet, consulté le 25 septembre 2026.
\end{omsources}

\section{Exercices}

\begin{exercise}[$\star$]\label{exo:s2:hedging-a-structured-products-book:1}
Une note vaut 99.12 pour 100 et la banque en vend pour 400~millions d'euros au pair. Quelle est la marge~?
\end{exercise}

\begin{exercise}[$\star$]\label{exo:s2:hedging-a-structured-products-book:2}
Le portefeuille est long de 9.66~millions d'euros de vega par point et le marché absorbe 5~millions d'euros par point. De combien la volatilité implicite baisse-t-elle, et
combien la banque paie-t-elle~?
\end{exercise}

\begin{exercise}[$\star$]\label{exo:s2:hedging-a-structured-products-book:3}
La couverture en delta détient 106~millions d'euros de l'indice à 100~\% et 334~millions d'euros à 55~\%. Quelle transaction la banque effectue-t-elle lorsque l'indice baisse entre ces deux niveaux~?
\end{exercise}

\begin{exercise}[$\star\star$]\label{exo:s2:hedging-a-structured-products-book:4}
Pourquoi l'émetteur d'un autocallable est-il long de volatilité~?
\end{exercise}

\begin{exercise}[$\star\star$]\label{exo:s2:hedging-a-structured-products-book:5}
Pourquoi l'émetteur d'une note worst-of est-il short de corrélation~?
\end{exercise}

\begin{exercise}[$\star\star$]\label{exo:s2:hedging-a-structured-products-book:6}
Pourquoi la couverture au voisinage des barrières peut-elle déstabiliser le marché~?
\end{exercise}

\begin{exercise}[$\star\star\star$]\label{exo:s2:hedging-a-structured-products-book:7}
\emph{Programmation.} Lire le profil de vega. Pour quel niveau de l'indice le vega d'une note isolée change-t-il de signe, et pourquoi~?
\end{exercise}

\begin{exercise}[$\star\star\star$]\label{exo:s2:hedging-a-structured-products-book:8}
\emph{Trouver la faille.} «~Nos notes rapportent une marge de 0.9~\% avec le risque couvert~; le recyclage n'est qu'un détail opérationnel.~»
\end{exercise}

\section{Problème~: long de ce que personne ne veut}

\begin{problem}[{Problème du week-end --- couverture d'un portefeuille de produits structurés}]\label{pb:s2:hedging-a-structured-products-book:1}
Le portefeuille synthétique du chapitre et les données publiques.

\medskip\noindent\textbf{Partie I --- Le portefeuille.}
\begin{enumerate}
\item Définir une \omterm{def:s2:hedging-a-structured-products-book:exposure}{exposition du portefeuille d'exotiques}.
\item Décrire les notes et le marché.
\item Quelle est la marge de la banque~?
\item Qu'a observé la BRI en Corée~?
\end{enumerate}

\medskip\noindent\textbf{Partie II --- Expositions.}
\begin{enumerate}[resume]
\item Donner le tableau des expositions.
\item Pourquoi la banque est-elle longue de volatilité, de skew et de dividendes~?
\item Pourquoi est-elle short de corrélation~?
\item Comment le vega varie-t-il avec le niveau de l'indice~?
\end{enumerate}

\medskip\noindent\textbf{Partie III --- Recyclage.}
\begin{enumerate}[resume]
\item Définir un \omterm{def:s2:hedging-a-structured-products-book:recycling}{trade de recyclage}.
\item Comment chaque exposition est-elle recyclée~?
\item Donner le tableau de recyclage.
\item Qui est payé, et par qui~?
\end{enumerate}

\medskip\noindent\textbf{Partie IV --- Le verdict.}
\begin{enumerate}[resume]
\item Énoncer le \emph{résultat nommé}~: les expositions nettes du portefeuille et le mouvement de volatilité implicite que provoque son recyclage.
\item Pourquoi la couverture est-elle déstabilisante en dessous de 55~\%~?
\item Comment le coût du recyclage se compare-t-il à la marge~?
\item Que se passe-t-il lorsque de nombreux émetteurs recyclent en même temps~?
\item Quels chapitres précédents se trouvent de l'autre côté de ce portefeuille~?
\item Quelle fiche de stratégie relève du trader et non de l'émetteur~?
\item Comment mener un backtest d'un acheteur de recyclage de vega~?
\item En une phrase~: que vend le desk de produits structurés en plus des notes~?
\end{enumerate}
\end{problem}

\section{Questions d'entretien}

\begin{interviewq}[$\star$ \iqroles{trader}]\label{iq:s2:hedging-a-structured-products-book:1}
Quelles expositions la banque conserve-t-elle après avoir vendu des autocallables~?
\end{interviewq}

\begin{interviewq}[$\star\star$ \iqroles{chercheur}]\label{iq:s2:hedging-a-structured-products-book:2}
Comment mesureriez-vous l'exposition à la corrélation d'un portefeuille de produits structurés~?
\end{interviewq}

\begin{interviewq}[$\star\star$ \iqroles{trader}]\label{iq:s2:hedging-a-structured-products-book:3}
L'indice a chuté de 40~\% et votre portefeuille est proche de ses barrières d'activation. Que faites-vous~?
\end{interviewq}

\begin{interviewq}[$\star\star$ \iqroles{risque}]\label{iq:s2:hedging-a-structured-products-book:4}
Quels tests de résistance comptent le plus pour un portefeuille d'autocallables~?
\end{interviewq}

\begin{interviewq}[$\star\star$ \iqroles{développeur}]\label{iq:s2:hedging-a-structured-products-book:5}
Réévaluer un portefeuille de dix mille notes pour vingt chocs prend trop de temps. Que feriez-vous~?
\end{interviewq}

\begin{interviewq}[$\star\star\star$ \iqroles{chercheur}]\label{iq:s2:hedging-a-structured-products-book:6}
Avec une concession de prix linéaire, montrer que la vente d'une quantité $V$ dans une profondeur $D$ coûte $V^2/(2D)$, et que le découpage de la vente en $n$ parts égales,
séparées par un rétablissement complet du prix, coûte $V^2/(2nD)$.
\end{interviewq}
